\section{Метод конечных разностей}
\label{sec:finite-difference-method}

\subsection{Основная идея}

Метод конечных разностей заменяет дифференциальную задачу системой алгебраических
уравнений на сетке. На равномерной одномерной сетке
\[
    x_i=a+ih,
    \qquad i=0,\ldots,N,
    \qquad h=\frac{b-a}{N}
\]
производные аппроксимируют разностями, полученными из формулы Тейлора:
\[
\begin{aligned}
 u'(x_i)&=\frac{u_{i+1}-u_i}{h}+O(h),\\
 u'(x_i)&=\frac{u_{i+1}-u_{i-1}}{2h}+O(h^2),\\
 u''(x_i)&=\frac{u_{i-1}-2u_i+u_{i+1}}{h^2}+O(h^2).
\end{aligned}
\]

Разностная схема должна обладать аппроксимацией, устойчивостью и сходимостью.
Локальная погрешность аппроксимации получается подстановкой точного решения в
схему. Устойчивость означает ограниченность роста возмущений. Для линейной
корректно поставленной начальной задачи теорема Лакса утверждает: при наличии
аппроксимации устойчивость эквивалентна сходимости.

\subsection{Уравнение теплопроводности}

Рассмотрим
\[
    u_t=a^2u_{xx}+f(x,t),
    \qquad 0<x<L,
    \qquad 0<t\leq T.
\]
Введём сетку $x_i=ih$, $t^n=n\tau$ и обозначим
$u_i^n\approx u(x_i,t^n)$.

Явная схема:
\[
    \frac{u_i^{n+1}-u_i^n}{\tau}
    =a^2\frac{u_{i-1}^n-2u_i^n+u_{i+1}^n}{h^2}+f_i^n.
\]
Она имеет порядок $O(\tau+h^2)$ и переписывается как
\[
    u_i^{n+1}
    =r u_{i-1}^n+(1-2r)u_i^n+r u_{i+1}^n+\tau f_i^n,
    \qquad
    r=\frac{a^2\tau}{h^2}.
\]
Условие устойчивости:
\[
    r\leq\frac12.
\]

Неявная схема:
\[
    \frac{u_i^{n+1}-u_i^n}{\tau}
    =a^2\frac{u_{i-1}^{n+1}-2u_i^{n+1}+u_{i+1}^{n+1}}{h^2}+f_i^{n+1}.
\]
Она безусловно устойчива, но на каждом шаге требует решения трёхдиагональной
системы. Порядок также $O(\tau+h^2)$.

Схема Кранка--Николсон усредняет пространственный оператор на слоях $n$ и $n+1$:
\[
\frac{u_i^{n+1}-u_i^n}{\tau}
=
\frac{a^2}{2h^2}
\left[
(u_{i-1}^{n}-2u_i^{n}+u_{i+1}^{n})+
(u_{i-1}^{n+1}-2u_i^{n+1}+u_{i+1}^{n+1})
\right]
+f_i^{n+1/2}.
\]
Она имеет порядок $O(\tau^2+h^2)$ и безусловно устойчива в энергетической норме.

\subsection{Волновое уравнение}

Для задачи
\[
    u_{tt}=a^2u_{xx}+f
\]
стандартная трёхслойная схема:
\[
    \frac{u_i^{n+1}-2u_i^n+u_i^{n-1}}{\tau^2}
    =a^2\frac{u_{i-1}^n-2u_i^n+u_{i+1}^n}{h^2}+f_i^n.
\]
Порядок аппроксимации --- $O(\tau^2+h^2)$. Условие Куранта:
\[
    \frac{a\tau}{h}\leq1.
\]
Для запуска схемы, кроме $u_i^0$, нужен первый временной слой $u_i^1$, который
строят по начальному условию для скорости и разложению Тейлора.

\subsection{Эллиптическое уравнение}

Для уравнения Пуассона
\[
    -\Delta u=f
\]
на прямоугольной сетке с шагами $h_x$, $h_y$ используется пятиточечный шаблон:
\[
-\frac{u_{i-1,j}-2u_{i,j}+u_{i+1,j}}{h_x^2}
-\frac{u_{i,j-1}-2u_{i,j}+u_{i,j+1}}{h_y^2}
=f_{i,j}.
\]
При $h_x=h_y=h$ схема имеет порядок $O(h^2)$. Получается большая разреженная
система. Для её решения применяют методы Зейделя, релаксации, сопряжённых
градиентов, многосеточные методы.

\subsection{Граничные условия}

Условие Дирихле задаётся непосредственным присваиванием узловых значений.
Условие Неймана
\[
    u_x(a)=q
\]
можно аппроксимировать односторонней разностью
\[
    \frac{u_1-u_0}{h}=q+O(h)
\]
или схемой второго порядка. Фиктивные узлы позволяют сохранить симметричный
шаблон. Для условия Робена $\alpha u+\beta u_n=g$ комбинируют значения функции
и разностную производную.

Важно, чтобы порядок аппроксимации граничных условий не снижал общий порядок
схемы.

\subsection{Исследование устойчивости}

Для линейных задач с постоянными коэффициентами применяется метод фон Неймана.
Погрешность представляют гармоникой
\[
    \varepsilon_i^n=G^n e^{\mathrm{i}k x_i},
\]
где $G$ --- коэффициент усиления. Условие устойчивости требует
\[
    |G|\leq1
\]
для всех допустимых волновых чисел.

Для нелинейных задач и переменных коэффициентов используют энергетические оценки,
принцип максимума и исследование спектра матрицы перехода.

\subsection{Преимущества и ограничения}

Преимущества: простота построения на регулярных сетках, прозрачность шаблонов,
эффективность для прямоугольных областей. Ограничения: трудности на сложной
геометрии, при локальном сгущении сетки и при необходимости точного соблюдения
законов сохранения. Для сложных областей часто предпочтительны методы конечных
объёмов или конечных элементов.

\subsection{Что важно сказать на экзамене}

Нужно показать переход от производных к сеточным разностям, определить
аппроксимацию, устойчивость и сходимость. В качестве примеров разумно привести
явную и неявную схемы теплопроводности, условие Куранта для волнового уравнения и
пятиточечный шаблон для уравнения Пуассона.

\newpage